\documentclass[./main]{subfiles}

% \graphicspath{{\subfix{../figures/section2/}}}


\begin{document}

\setcounter{section}{1}

\section{\TikZ}

\subsection{What}
\begin{frame}[fragile]
  \frametitle{\subsecname\ \secname ?}

  \centering

  \TikZ ist \textit{kein} Zeichenprogramm (German) \\
  \TikZ は描画プログラムではない{\small \citep{WikiTikZEn,tikz-manual}}

  \vspace{1em}

  テキストベースで記述された命令を元に図形を生成する

  \vspace{1em}

  \begin{minipage}{0.8\linewidth}
    \begin{verbatim}
\usepackage{tikz}
\begin{tikzpicture}
  \draw(0,0)--(1,0)--(1,1)--cycle;
\end{tikzpicture}
    \end{verbatim}  
  \end{minipage}
  \begin{minipage}{0.15\linewidth}
    \begin{tikzpicture}
      \draw(0,0)--(1,0)--(1,1)--cycle;
    \end{tikzpicture}
  \end{minipage}

\end{frame}


\subsection{Why}
\begin{frame}
  \frametitle{\subsecname\ \secname ?}

  \centering

  \vfilll

  \begin{itemize}
    \item ベクトルベースなので図形が解像度に依存しない
    \item 文書との一貫性
    \item 豊富なライブラリ
    \item テキストエディタのみで編集可能
    \item 共有・修正が楽
  \end{itemize}

  \vfilll

  \begin{itemize}
    \item 使い道はさておき，パワポではできなさそうな\\ 例をいくつか挙げる
  \end{itemize}

  \vfilll

\end{frame}





\subsection{Application}

\newcommand{\drawClock}[2]{%
  \begin{tikzpicture}
    \draw[very thick] (0,0) circle (2cm);%時計の外周
    \draw[fill] (0,0) circle (0.1cm);%中心点
    \foreach \angle / \label in {0/3, 30/2, 60/1, 90/12, 120/11, 150/10, 180/9, 210/8, 240/7, 270/6, 300/5, 330/4} { 
      \draw (\angle:1.8cm) -- (\angle:2cm); 
      \node at (\angle:1.5cm) {\textsf{\label}}; 
    } 
    \foreach \angle in {0,90,180,270} 
       \draw[thick] (\angle:1.7cm) -- (\angle:2cm); 
    \draw[line width=3pt, cap=round] (0,0) -- ($(-#1*30-#2*0.5+90:0.8cm)$);%短針
    \draw[line width=3pt, cap=round] (0,0) -- ($(-#2*6+90:1.2cm)$);%長針
  \end{tikzpicture}%
}



\begin{frame}[allowframebreaks,fragile]
  \frametitle{\secname\ \subsecname}

  \centering

  \begin{figure}[htbp]
    \centering
    \begin{minipage}{0.45\linewidth}
      \centering
      \reinitrand[counter=random,first=0,last=23]  \rand
      \reinitrand[counter=random2,first=0,last=59] \rand
      \drawClock{\arabic{random}}{\arabic{random2}}
      \caption{1分ごとにランダム}
    \end{minipage}
    \begin{minipage}{0.45\linewidth}
      \centering
      \DTMsavenow{now}
      \drawClock{\DTMfetchhour{now}}{\DTMfetchminute{now}}
      \caption{コンパイル時刻}
    \end{minipage}
  \end{figure}

  \framebreak

  \begin{minipage}{0.45\linewidth}
    \begin{figure}[htbp]
      \centering
      \includestandalone[width=\linewidth]{TopologicalDefect}
      \caption{\centering Topological defect $(k,c)$\\\citep{topo,tikz-topo}}
    \end{figure}
  \end{minipage}
  \hfill
  \begin{minipage}{0.5\linewidth}
    閉曲線$\Gamma$,\\
    棒の角度$\theta(\bm{r})=k\phi(\bm{r})+c$,\\
    方位角$\phi(\bm{r})=\tan^{-1}(x/y)$
    \begin{equation*}
      \oint_\Gamma\diff\theta=\oint_\Gamma\frac{\diff\theta}{\diff l}\diff l=2\pi k
    \end{equation*}
  \end{minipage}  

  \framebreak
  \begin{figure}[htbp]
    \centering
    \includestandalone[width=0.5\linewidth]{TriangleIteration}
    \caption{Turtle graphics}
  \end{figure}

  \framebreak
  \begin{figure}[htbp]
    \centering
    \includestandalone[width=\linewidth]{CameraAngle}
    \caption{視点を少し変える\citep{tikz-3dplot}}
  \end{figure}

  \framebreak
  \centering
  自作関数一つで色々なニューラルネットが書ける\\


  \newcommand{\drawNN}[5][1]{
  \begin{tikzpicture}[scale=#1]


  \newcommand{\numhh}{#2}
  \newcommand{\numi}{#3}
  \newcommand{\numh}{#4}
  \newcommand{\numo}{#5}

  \newcommand{\dist}{1.5}
  
  \draw [decorate,decoration={brace,amplitude=10},xshift=0pt,yshift=0pt](\dist-\dist/2,\numh/2+0.5) -- (\dist*\numhh+\dist/2,\numh/2+0.5) node [above right of=1,black,midway,xshift=-0.6cm]{hidden layer(s)};
  \draw [decorate,decoration={brace,amplitude=10},xshift=0pt,yshift=0pt](-0.2,-\numi/2-0.5) -- (-0.2,\numi/2+0.5) node [left=0,black,midway,xshift=-0.6cm,align=right]{input\\layer};
  \draw [decorate,decoration={brace,mirror,amplitude=10},xshift=0pt,yshift=0pt](\dist*\numhh+\dist*1.2,-\numo/2) -- (\dist*\numhh+\dist*1.2,\numo/2) node [right=1cm,black,midway,xshift=-0.6cm,align=left]{output\\layer};

  \tikzstyle{input}=[draw,fill=red!50,circle,minimum size=10pt,inner sep=0pt]
  \tikzstyle{hidden}=[draw,fill=green!50,circle,minimum size=10pt,inner sep=0pt]
  \tikzstyle{output}=[draw,fill=blue!50,circle,minimum size=10pt,inner sep=0pt]
  \tikzstyle{bias}=[draw,fill=gray!50,circle,minimum size=10pt,inner sep=0pt]
  \tikzstyle{arrow}=[arrows={{Latex[scale=0.5]}-}, thick]  % , in=90, out=-90


  \pgfmathsetmacro{\y}{(\numi)/2}
  \node (n00)[bias] at (0,\y) {};
  \foreach \i in {1, ..., \numi}{
    \pgfmathsetmacro{\y}{-(\numi+2)/2+\i}
    \node (n0\i)[input] at (0,\y) {};
  }


  \foreach \j in {1, ..., \numhh}{
    \pgfmathsetmacro{\x}{\dist*\j}
    \pgfmathsetmacro{\y}{(\numh)/2}
    \node (n\j0)[bias] at (\x,\y) {};

    \foreach \i in {1, ..., \numh}{
      \pgfmathsetmacro{\y}{-(\numh+2)/2+\i}
      \node (n\j\i)[hidden] at (\x,\y) {};
    }
  }

  \pgfmathsetmacro{\j}{\numhh+1}
  \pgfmathsetmacro{\x}{\dist*\j}
  \foreach \i in {1, ..., \numo}{
    \pgfmathsetmacro{\y}{-(\numo+1)/2+\i}
    \node (no\i)[output] at (\x,\y) {};
  }

  \foreach \j in {1, ..., \numhh}{
    \pgfmathsetmacro{\jj}{int(\j-1)}
    \foreach \i in {1, ..., \numh}{
      \ifnum\j=1
        \foreach \k in {0, ..., \numi}{
          \draw (n\j\i) edge[arrow] (n0\k);
        }
      \else
        \foreach \k in {0, ..., \numh}{
          \draw (n\j\i) edge[arrow] (n\jj\k);
        }
      \fi
    }
  }

  \pgfmathsetmacro{\j}{\numhh}
  \foreach \i in {1, ..., \numo}{
    \foreach \k in {0, ..., \numh}{
      \draw (no\i) edge[arrow] (n\j\k);
    }
  }
\end{tikzpicture}
}
  \begin{minipage}[b]{0.35\linewidth}
    \centering
    \scalebox{0.6}{\drawNN{1}{3}{2}{1}}
    \verb|\drawNN{1}{3}{2}{1}|
  \end{minipage}
  \hfill
  \begin{minipage}[b]{0.6\linewidth}
    \centering
    \scalebox{0.6}{\drawNN{3}{7}{5}{3}}
    \verb|\drawNN{3}{7}{5}{3}|
  \end{minipage}
  
\end{frame}




\ifSubfilesClassLoaded{%
  \begin{frame}
    \bibliography{../ref}
  \end{frame}
}{}


\end{document}